\section{Discussion}
In this section, we reflect on the applications, future work, limitations, and ethical and societal risks of generative agents.

\subsection{Applications of Generative Agents}
Generative agents have vast potential applications that extend beyond the sandbox demonstration presented in this work, especially in domains that would benefit from a model of human behavior based on long-term experience. For instance, social simulacra have demonstrated the ability to create stateless personas that generate conversation threads in online forums for social prototyping~\cite{9_park2022socialsimulacra}. \crdraft{With generative agents, we can populate these forums, as well as virtual reality metaverses~\cite{orland2021metaverse} or   physical spaces with social robots~\cite{12_bartneck2004design} if paired with multimodal models. This opens up the possibility of creating even more powerful simulations of human behavior to test and prototype social systems and theories, as well as to create new interactive experiences.}

Another application area is in the human-centered design process, similar to the intended applications of cognitive models such as GOMS~\cite{john1996goms} and the KLM~\cite{card1980keystroke}. Consider a generative agent that models Sal, the protagonist in Mark Weiser's famous ubiquitous computing vignette~\cite{weiser1991computer}, based on her life patterns and interactions with technology. In this scenario, the agent acts as a proxy for Sal and learns plausible sets of behaviors and reflections that Sal may exhibit based on her life. The agent can encode information such as when Sal wakes up, when she needs her first cup of coffee, and what her typical day looks like. Using this information, the agent can automatically brew coffee, help get the kids ready for school, and adjust the ambient music and lighting to match Sal's mood after a hard day at work. By utilizing generative agents as proxies for users, we can develop a deeper understanding of their needs and preferences, resulting in more personalized and effective technological experiences.

\subsection{Future Work and Limitations}
In this work, we introduced generative agents and presented an initial implementation and evaluation of their architecture. Future research can build upon the proposed agent architecture to improve and further evaluate its performance. In terms of implementation, the retrieval module, for example, could be enhanced to retrieve more relevant information given a context by fine-tuning the relevance, recency, and importance functions that compose the retrieval function. Additionally, efforts can be made to improve the architecture's performance, making it more cost-effective. The present study required substantial time and resources to simulate 25 agents for two days, costing thousands of dollars in token credits and taking multiple days to complete. To enhance real-time interactivity, future work can explore parallelizing agents or developing language models specifically designed for building generative agents. In general, with advances in underlying models, we believe that agents' performance will improve.

In terms of evaluation, the assessment of generative agents' behavior in this study was \crdraft{limited to a relatively short timescale and a baseline human crowdworker condition. While the crowdworker condition provided a helpful comparison point, it did not represent the maximal human performance that could serve as the gold standard in terms of believability. Future research should aim to observe the behavior of generative agents over an extended period to gain a more comprehensive understanding of their capabilities and establish rigorous benchmarks for more effective performance testing.} Additionally, varying and contrasting the underlying models, as well as the hyperparameters used for the agents during future simulations, could provide valuable insights into the impact of these factors on the agents' behavior. Lastly, the robustness of generative agents is still largely unknown. They may be vulnerable to prompt hacking, memory hacking---where a carefully crafted conversation could convince an agent of the existence of a past event that never occurred---and hallucination, among other issues. Future research can comprehensively test these robustness concerns, and as large language models become more resilient to such attacks, generative agents can adopt similar mitigations.

In general, any imperfections in the underlying large language models will be inherited by generative agents. Given the known biases of language models, generative agents may potentially exhibit biased behavior or stereotypes. Moreover, like many large language models, generative agents may struggle to generate believable behavior for certain subpopulations, particularly marginalized populations, due to limited data availability. While improvements to the agents' modules may mitigate some of these issues, we believe that addressing them fundamentally requires improving the underlying large language models by aligning their values with the desired outcomes of the agents.

\subsection{Ethics and Societal Impact}
Generative agents, while offering new possibilities for human-computer interaction, also raise important ethical concerns that must be addressed. \crdraft{One risk is people forming parasocial relationships with generative agents, even when such relationships may not be appropriate. Despite being aware that generative agents are computational entities, users may anthropomorphize them or attach human emotions to them~\cite{hofstadter1995fluid, reeves1996media}. While this tendency may increase user engagement, it also poses risks, such as users becoming overly reliant on or emotionally attached to the agents~\cite{abercrombie2023mirages}. To mitigate this risk, we propose two principles. First, generative agents should explicitly disclose their nature as computational entities. Second, developers of generative agents must ensure that the agents, or the underlying language models, are value-aligned so that they do not engage in behaviors that would be inappropriate given the context, for example, reciprocating confessions of love.}

A second risk is the impact of errors. For example, if a ubiquitous computing application makes the wrong inference about a user's goals based on generative agent predictions, it could lead to annoyance at best and outright harm at worst. In our instantiation of generative agents, we mitigate these risks by focusing on an interactive video game environment, where such harms are unlikely. However, in other application domains, it will be important to follow best practices in human-AI design~\cite{amershi2019guidelines, yang2020re} to understand errors and how they might percolate into the user experience.

Third, generative agents may exacerbate existing risks associated with generative AI, such as deepfakes, misinformation generation, and tailored persuasion. To mitigate this risk, we suggest that platforms hosting generative agents maintain an audit log of the inputs and generated outputs. This would enable the detection, verification, and intervention against malicious use. While logging alone cannot directly prevent such misuse, it can reduce the likelihood of motivated actors engaging in this behavior, as the risk of disclosure would be higher. Additionally, building this architecture oneself can be time-consuming (in our case, roughly a year), which may deter some actors from pursuing such behavior by using their own generative agent infrastructures.

A fourth risk is over-reliance: the concern that developers or designers might use generative agents and displace the role of humans and system stakeholders in the design process~\cite{9_park2022socialsimulacra}. We suggest that generative agents should never be a substitute for real human input in studies and design processes. Instead, they should be used to prototype ideas in the early stages of design when gathering participants may be challenging or when testing theories that are difficult or risky to test with real human participants. By adhering to these principles, we can ensure that the deployment of generative agents in the wild is ethical and socially responsible.
